% !TeX root = ../教材题库_全部条目_转写工作稿.tex
% Source: XJ-B1-C05-HW-P162-Y0217
% 第5章 / 5.1 任意角和弧度制 / 教材 p.162 / PDF 第169页
% OCR未校对：公式与文字均可能有误；未完成数学TeX转写；禁止直接发给学生
% 原题及相邻原解截图（核对用）：../题目原图/XJ-B1-C05-HW-P162-Y0217-01.webp
% 本文件只在转写完成后作为题干使用；OCR 条目尚不是可靠题干。
\begin{bankitem}{习题5.1 第7题}{第5章 / 5.1 任意角和弧度制 / 教材 p.162 / PDF 第169页}{XJ-B1-C05-HW-P162-Y0217}{OCR未校对：公式与文字均可能有误；未完成数学TeX转写；禁止直接发给学生}
\begin{ocrdraft}
[教材p.162 / PDF第169页]
\par
7．在航海上，最初规定地球子午圈的1分弧长为1海里（如图），已知地球的半
\par
径约为6371km，根据此规定计算1海里的长度．
\par
北极点
\par
1海里
\par
1'囱心角
\par
赤道
\par
南极点
\par
（第7题）
\par
（第8题）
\end{ocrdraft}
\end{bankitem}
